

\section{问题三的建模与求解}

\subsection{光伏发电功率预报处理}

问题三在问题二基础上引入$0{:}00$、$6{:}00$、$12{:}00$和$18{:}00$
发布的未来24 h光伏发电功率预报。每天$0{:}00$根据当时获得的预报制定
全天计划，后续预报时刻仅调整尚未执行时段的购电量。\subsection{模型基本设定}

问题三沿用前两问的时间尺度、功率与电量换算方式及储能参数，每天划分为144个$10\,\mathrm{min}$时段，$\Delta t=1/6\,\mathrm{h}$。负载预测采用问题二的Dynamic模型，并根据$0{:}00$、$6{:}00$、$12{:}00$和$18{:}00$获得的光伏发电功率预报调整尚未执行时段的购电量。计算从2025年1月1日开始，以1月数据形成历史信息并传递储电量状态，正式统计期为2月1日至12月31日，共334天；储电量按
\begin{equation}
E_{d+1,0}=E_{d,144},\qquad E_{2025\text{-}01\text{-}01,0}=6000\,\mathrm{kWh}
\end{equation}
进行跨日传递。





附件3提供未来24 h整点光伏发电功率预报，而购电策略采用
$10\,\mathrm{min}$调度尺度。为使两者对应，对相邻整点预报进行线性插值。
设发布时间为$\tau$，相邻整点$h$和$h+1$的预测功率分别为
$P_{d,h}^{f(\tau)}$和$P_{d,h+1}^{f(\tau)}$，则区间内时刻
$s\in[h,h+1]$的预测功率为
\begin{equation}
P_{d}^{f(\tau)}(s)
=
P_{d,h}^{f(\tau)}
+
(s-h)
\left[
P_{d,h+1}^{f(\tau)}
-
P_{d,h}^{f(\tau)}
\right].
\label{eq:p3_pv_interp}
\end{equation}

取各$10\,\mathrm{min}$时段中点的插值功率，并乘以时段长度$\Delta t$，
得到对应的光伏发电预测电量：
\begin{equation}
R_{d,t}^{f(\tau)}
=
P_{d}^{f(\tau)}(s_t)\Delta t.
\label{eq:p3_pv_energy}
\end{equation}
其中，$s_t$表示第$t$个时段的中点；发布时间至第一个预测整点之间的时段，
采用第一个整点预测功率补齐。

典型日不同发布时间的光伏预测如图~\ref{fig:06}所示。
随着发布时间推移，可用于后续时段决策的光伏预测不断更新，
为日内调整购电量提供新的信息。

\begin{figure}[H]
    \centering
    \includegraphics[width=1\linewidth]{图06_典型日光伏预测逐批更新.png}
    \caption{典型日光伏预测逐批更新}
    \label{fig:06}
\end{figure}


\subsection{计划购电模型}

每天$0{:}00$根据当日负载预测和光伏发电功率预报制定计划购电策略。
设第$d$天第$t$个时段的计划购电量为$x_{d,t}^{\mathrm{plan}}$，
储能充电量和放电量分别为$u_{d,t}$和$v_{d,t}$，
储电量为$E_{d,t}$，弃光量为$w_{d,t}$。
各时段满足
\begin{equation}
x_{d,t}^{\mathrm{plan}}
+R_{d,t}^{f(0)}
+v_{d,t}
+\widehat e_{d,t}^{(0)}
=
\widetilde L_{d,t}
+u_{d,t}
+w_{d,t}.
\label{eq:p3_plan_balance}
\end{equation}

其中，$R_{d,t}^{f(0)}$为$0{:}00$预报对应的光伏发电量，
$\widetilde L_{d,t}$为当日负载预测量，
$\widehat e_{d,t}^{(0)}$为预测供电缺口对应的紧急购电量。
以计划购电费用和预测紧急购电费用之和最小为目标：
\begin{equation}
\min C_d^{\mathrm{plan}}
=
\sum_{t=1}^{144}
\left(
c_t x_{d,t}^{\mathrm{plan}}
+
5c_t\widehat e_{d,t}^{(0)}
\right).
\label{eq:p3_plan_obj}
\end{equation}

储电量递推、容量限制、最大充放电功率及充放电互斥约束均沿用第4.4节。
求解后得到$0{:}00$确定的全天计划购电量及储能充放电方案。


\subsection{日内调整模型}

在$6{:}00$、$12{:}00$和$18{:}00$获得新的光伏发电功率预报后，
仅重新计算尚未执行时段的购电量。设当前调整时刻为$\tau$，
其后尚未执行的时段集合为$\mathcal{T}_{\tau}$。
已经执行的购电量保持不变，负载预测继续采用$0{:}00$确定的结果，
仅更新光伏发电功率预报。

考虑紧急购电价格为正常电价的5倍，本文采用较为保守的调整规则：
保持$0{:}00$确定的储能充放电方案不变，仅允许在原计划基础上增加购电量。
设$x_{d,t}^{(\tau)}$为调整后的购电量，
$d_{d,t}^{+(\tau)}$为相对原计划的增加量，则
\begin{equation}
\left\{
\begin{aligned}
x_{d,t}^{(\tau)}
&=x_{d,t}^{\mathrm{plan}}+d_{d,t}^{+(\tau)},\\
d_{d,t}^{+(\tau)}
&\geq0,
\end{aligned}
\right.
\qquad t\in\mathcal{T}_{\tau}.
\label{eq:p3_adjust_constraint}
\end{equation}

根据题目规定，增加购电量按交易时刻电价的1.5倍计价。
因此，每次调整以剩余时段的计划购电费用、增加购电费用和预测紧急购电费用
之和最小为目标：
\begin{equation}
\min C_d^{(\tau)}
=
\sum_{t\in\mathcal{T}_{\tau}}
\left[
c_t x_{d,t}^{\mathrm{plan}}
+
1.5c_t d_{d,t}^{+(\tau)}
+
5c_t\widehat e_{d,t}^{(\tau)}
\right].
\label{eq:p3_adjust_obj}
\end{equation}

其中，$\widehat e_{d,t}^{(\tau)}$为根据时刻$\tau$最新光伏预报得到的
预测紧急购电量。每次获得新预报后，分别计算沿用上一轮已接受方案和候选调整方案
的预计费用；仅当候选方案费用更低时才更新尚未执行时段的购电量。
所有调整量及相关费用均以$0{:}00$原计划为基准计算。


\subsection{购电策略与结论分析}



问题三采用日内滚动调整策略。每天$0{:}00$根据Dynamic负载预测和当时获得的光伏发电功率预报，确定全天计划购电量及储能充放电方案；在$6{:}00$、$12{:}00$和$18{:}00$获得新的光伏发电功率预报后，保持已执行购电量和$0{:}00$确定的储能充放电方案不变，\textbf{仅允许增加尚未执行时段的购电量}。每次分别计算沿用当前方案与候选调整方案的预计费用，\textbf{仅当候选方案费用更低时接受调整，否则继续执行当前方案}。




根据统计期内已有计算结果，相较仅执行$0{:}00$计划的方案，利用$6{:}00$、$12{:}00$和$18{:}00$新增预报进行调整，使紧急购电量下降约$60.94\%$，总购电费用减少约$12.07$万元。因此，整体上有必要利用日内新增光伏预报制定调整购电策略。但部分月份和日期的实际总费用仍有所增加，说明预报信息更新并不保证每次调整都能在实际结算中节省费用，应按前述预计费用比较规则决定是否接受调整。

\subsection{指定日期的购电与储能结果}

按照题面表1、表2和表3的要求，下面列出四个指定日期的购电、储能和紧急购电结果。表~\ref{tab:p3_purchase}中的六个指定时段购电量为$0{:}00$确定的原计划，全天调整后购电量为日内预报更新后实际执行的非紧急购电量；全天总购电费用包含计划、调整及紧急购电费用。

\begin{table}[H]
\centering
\caption{问题三指定日期的计划购电量、全天购电量与费用}
\label{tab:p3_purchase}
\resizebox{\textwidth}{!}{
\begin{tabular}{cccccccc}
\toprule
\multicolumn{2}{c}{2025-03-20} & \multicolumn{2}{c}{2025-06-21} & \multicolumn{2}{c}{2025-09-23} & \multicolumn{2}{c}{2025-12-21} \\
\cmidrule(lr){1-2}\cmidrule(lr){3-4}\cmidrule(lr){5-6}\cmidrule(lr){7-8}
时间段 & 购电量（kWh） & 时间段 & 购电量（kWh） & 时间段 & 购电量（kWh） & 时间段 & 购电量（kWh） \\
\midrule
10:00--10:10 & 0.0000 & 10:00--10:10 & 0.0000 & 10:00--10:10 & 0.0000 & 10:00--10:10 & 0.0000 \\
12:00--12:10 & 539.2991 & 12:00--12:10 & 0.0000 & 12:00--12:10 & 535.3132 & 12:00--12:10 & 985.0206 \\
14:00--14:10 & 0.0000 & 14:00--14:10 & 0.0000 & 14:00--14:10 & 0.0000 & 14:00--14:10 & 0.0000 \\
16:00--16:10 & 598.3578 & 16:00--16:10 & 195.9234 & 16:00--16:10 & 655.5594 & 16:00--16:10 & 828.9816 \\
18:00--18:10 & 731.3515 & 18:00--18:10 & 0.0000 & 18:00--18:10 & 852.6645 & 18:00--18:10 & 736.7735 \\
20:00--20:10 & 0.0000 & 20:00--20:10 & 0.0000 & 20:00--20:10 & 30.6555 & 20:00--20:10 & 0.0000 \\
\midrule
全天计划购电量（kWh） & 68484.9836 & 全天计划购电量（kWh） & 34535.9755 & 全天计划购电量（kWh） & 81380.3581 & 全天计划购电量（kWh） & 99015.9587 \\
全天计划购电费用（元） & 41528.8937 & 全天计划购电费用（元） & 19529.8668 & 全天计划购电费用（元） & 51350.5106 & 全天计划购电费用（元） & 63706.4610 \\
全天调整后购电量（kWh） & 69714.5663 & 全天调整后购电量（kWh） & 36722.0850 & 全天调整后购电量（kWh） & 81644.4435 & 全天调整后购电量（kWh） & 100348.7128 \\
全天总购电费用（元） & 50754.4824 & 全天总购电费用（元） & 24365.1911 & 全天总购电费用（元） & 51665.0276 & 全天总购电费用（元） & 65931.2930 \\
\bottomrule
\end{tabular}}
\end{table}

由表~\ref{tab:p3_purchase}可得，四天的全天计划购电量分别为$68,484.984$、$34,535.976$、$81,380.358$和$99,015.959$ kWh。6月21日计划购电量最低，12月21日最高，说明在不同日期的负载预测、光伏预报和储能状态下，外网购电需求存在明显差异。日内调整使四天非紧急购电量分别增加$1,229.583$、$2,186.110$、$264.085$和$1,332.754$ kWh；其中6月21日增购量最大，而9月23日增购量较小。四天最终总费用分别为$50,754.482$、$24,365.191$、$51,665.028$和$65,931.293$元。计划费用与最终总费用的差额同时包含调整费用和紧急购电费用，不能直接作为日内调整的节省额。

\begin{table}[H]
\centering
\caption{问题三指定日期最终储能充放电量及日初、日末储电量}
\label{tab:p3_storage}
\resizebox{\textwidth}{!}{
\begin{tabular}{cccccccccccc}
\toprule
\multicolumn{3}{c}{2025-03-20} & \multicolumn{3}{c}{2025-06-21} & \multicolumn{3}{c}{2025-09-23} & \multicolumn{3}{c}{2025-12-21} \\
\cmidrule(lr){1-3}\cmidrule(lr){4-6}\cmidrule(lr){7-9}\cmidrule(lr){10-12}
时间段 & 充电量（kWh） & 放电量（kWh） & 时间段 & 充电量（kWh） & 放电量（kWh） & 时间段 & 充电量（kWh） & 放电量（kWh） & 时间段 & 充电量（kWh） & 放电量（kWh） \\
\midrule
0:00--4:00 & 9000.0000 & 0.0000 & 0:00--4:00 & 1466.5557 & 0.0000 & 0:00--4:00 & 9000.0000 & 0.0000 & 0:00--4:00 & 9000.0000 & 0.0000 \\
4:00--8:00 & 1666.6667 & 6783.1390 & 4:00--8:00 & 833.3333 & 1862.9101 & 4:00--8:00 & 1666.6667 & 5034.1190 & 4:00--8:00 & 2500.0000 & 1508.3333 \\
8:00--12:00 & 7130.5153 & 1856.8610 & 8:00--12:00 & 817.6541 & 0.0000 & 8:00--12:00 & 5833.3333 & 3605.8810 & 8:00--12:00 & 5833.3333 & 7806.6667 \\
12:00--16:00 & 4745.7974 & 979.8133 & 12:00--16:00 & 9849.0125 & 0.0000 & 12:00--16:00 & 5952.5628 & 906.5759 & 12:00--16:00 & 7719.5487 & 2337.8344 \\
16:00--20:00 & 0.0000 & 5627.6462 & 16:00--20:00 & 0.0000 & 6008.6575 & 16:00--20:00 & 0.0000 & 5833.3333 & 16:00--20:00 & 0.0000 & 5559.6742 \\
20:00--24:00 & 0.0000 & 3012.3538 & 20:00--24:00 & 0.0000 & 2631.3425 & 20:00--24:00 & 0.0000 & 2806.6667 & 20:00--24:00 & 0.0000 & 3080.3258 \\
\midrule
0:00储电量 & 1200.0000 & -- & 0:00储电量 & 1200.0000 & -- & 0:00储电量 & 1200.0000 & -- & 0:00储电量 & 1200.0000 & -- \\
24:00储电量 & 1200.0000 & -- & 24:00储电量 & 1200.0000 & -- & 24:00储电量 & 1200.0000 & -- & 24:00储电量 & 1200.0000 & -- \\
\bottomrule
\end{tabular}}
\end{table}

表~\ref{tab:p3_storage}按六个4 h区间汇总最终执行的储能充、放电量。由于本问日内调整冻结$0{:}00$储能计划，表中最终储能调度与日初计划一致。四天的日初、日末储电量均为$1200.000$ kWh，处于储能容量下界；这一结果来自跨日连续运行的调度，而非强制每日恢复日初储电量。从区间分布看，3月20日、9月23日和12月21日均在0:00--4:00充电$9000.000$ kWh，而6月21日该区间仅充电$1466.556$ kWh，主要充电量集中在12:00--16:00，为$9849.013$ kWh。四天16:00--24:00均未安排充电，两个区间合计放电量均为$8640.000$ kWh，体现出前段补充储电、后段释放储能的调度特征。同一4 h区间内同时出现充电量和放电量，表示区间内不同时段分别充、放电，并不表示同一10 min时段同时充放电。

为突出实际发生供电缺口的时段，表~\ref{tab:p3_emergency}按日期列出最先发生的若干段非零紧急购电记录；连续发生紧急购电的10 min时段合并展示，其购电量为对应时段电量之和。表末总量按全天所有时段汇总，完整非零时段明细见附录（待补充）。

\begin{table}[H]
\centering
\caption{问题三指定日期的部分紧急购电时段及全天总量}
\label{tab:p3_emergency}
\resizebox{\textwidth}{!}{
\begin{tabular}{cccccccc}
\toprule
\multicolumn{2}{c}{2025-03-20} & \multicolumn{2}{c}{2025-06-21} & \multicolumn{2}{c}{2025-09-23} & \multicolumn{2}{c}{2025-12-21} \\
\cmidrule(lr){1-2}\cmidrule(lr){3-4}\cmidrule(lr){5-6}\cmidrule(lr){7-8}
时间段 & 购电量（kWh） & 时间段 & 购电量（kWh） & 时间段 & 购电量（kWh） & 时间段 & 购电量（kWh） \\
\midrule
1:00--1:30 & 42.1548 & 0:20--0:30 & 34.9023 & 20:50--21:00 & 12.5142 & 1:10--1:20 & 3.7376 \\
3:10--3:30 & 27.1762 & 1:20--1:40 & 46.1316 & -- & -- & 3:30--3:50 & 107.4211 \\
5:00--5:10 & 2.1611 & 1:50--2:00 & 13.1523 & -- & -- & 6:10--6:30 & 28.2280 \\
6:30--7:00 & 64.5172 & 2:50--3:30 & 99.7451 & -- & -- & 11:00--11:20 & 26.2342 \\
7:10--7:40 & 40.7259 & 4:00--4:10 & 11.6759 & -- & -- & 18:50--19:00 & 10.2994 \\
8:00--12:00 & 1661.9423 & 6:50--7:10 & 14.4389 & -- & -- & 19:50--20:00 & 27.7264 \\
\midrule
全天紧急购电量（kWh） & 1958.5509 & 全天紧急购电量（kWh） & 488.9160 & 全天紧急购电量（kWh） & 12.5142 & 全天紧急购电量（kWh） & 203.6468 \\
\bottomrule
\end{tabular}}
\end{table}

四天紧急购电总量分别为$1,958.551$、$488.916$、$12.514$和$203.647$ kWh。3月20日紧急购电量最高，9月23日最低；即使日内允许增加购电量，实际负载及光伏出力与预测之间仍可能存在偏差，因而紧急购电不能完全消除。紧急购电按交易时刻电价的5倍结算，风险评价需同时考虑缺口电量及其发生时段，不能仅以全天购电总量判断方案优劣。
